% =============================================================================
% Partie C — Conclusion
% =============================================================================
\section{Conclusion}

\subsection{Ce qui est tenu, et comment on le sait}

Deux contraintes étaient éliminatoires. Elles sont traitées par des mécanismes
différents, et vérifiées par des preuves différentes — ce qui est précisément
pourquoi elles méritent d'être distinguées.

\textbf{C1 est prouvée en charge, en production.} Un lien de 500 exemplaires a
reçu 16\,753 requêtes dont 500 réservations exactement, sans une seule erreur
serveur, avec un pic à 200 requêtes par seconde. Le compteur final est à zéro et
l'état du lien est « épuisé ». Ce résultat n'est pas une moyenne : la contrainte
est « exactement $N$ », et l'écart mesuré est nul.

\textbf{C2 est prouvée en deux temps.} Sous charge, chaque modification émise a
reçu son acquittement : aucun message manquant. Mais un acquittement ne prouve
pas la convergence du contenu, et un outil de charge ne peut pas le faire. Un
second outil, avec un vrai client CRDT, reconstruit le document depuis le journal,
vérifie que les numéros de séquence sont contigus, et compare deux
reconstructions indépendantes. Il conclut à la convergence, et il conclut à
l'échec dès qu'un numéro manque.

\textbf{Les quatre autres contraintes sont couvertes par des tests unitaires.}
Le quota de crédit, la version figée du PDF, le comportement sous pic et la
notification du dernier exemplaire ont chacun leur jeu de tests, sans aucun appel
à un service distant.

\subsection{Ce qui reste à mesurer}

Le dossier doit être honnête sur ce point.

Le nombre de rédacteurs simultanés effectivement tenu reste à établir par paliers
successifs, et le résultat doit être rapporté tel quel, même inférieur à la cible
de 300. Une cible non atteinte et annoncée vaut mieux qu'un chiffre flatteur et
invérifiable.

Le comportement après un arrêt brutal du processus en pleine écriture dispose
d'un scénario dédié, qui tue le processus et vérifie que le journal a progressé
sans trou après redémarrage. Ce scénario est écrit, son exécution reste à
consigner.

Les métriques système — occupation processeur, mémoire de pointe, solde de
crédits de processeur — ne sont pas observables depuis l'outil de charge. Elles
se relèvent en parallèle sur la supervision du fournisseur, et doivent
accompagner les résultats de charge pour que le dimensionnement soit justifié
plutôt qu'affirmé.

\subsection{Ce que ce projet a réellement appris}

Trois leçons, dans l'ordre de ce qu'elles ont coûté.

\textbf{Première leçon : un faux service doit imiter les refus, pas seulement les
succès.} Un double en mémoire de la base évaluait lui-même les expressions de
condition. Il était donc plus permissif que le vrai service. Résultat :
quarante-deux tests verts, et aucune écriture possible en production, parce que le
service réel refuse l'arithmétique entre deux valeurs dans une condition. Le même
défaut s'est reproduit avec les mots réservés. La correction n'a pas été de
réparer le code, mais de rendre le double \emph{plus strict} : il rejette
désormais ce que le service rejette.

\textbf{Deuxième leçon : une compilation réussie ne prouve rien.} Les douze
diagrammes de conception compilaient sans erreur, et plusieurs étaient faux : un
code 200 là où le service répond 202, une redirection qui n'existe pas, un
numéro de version compté deux fois. L'outil de diagrammes écrit ses erreurs dans
l'image et renvoie un succès. Il a fallu relire chaque source contre le code et
regarder chaque rendu.

\textbf{Troisième leçon : un test qui ne peut pas échouer ne protège rien.} Un
seuil de taux d'échec posé sur un parcours qui renvoie légitimement des refus est
inatteignable par construction. Un outil de vérification qui interroge une route
dont la réponse n'a pas la forme attendue ne vérifie rien. À chaque fois, la
question utile n'est pas « le test passe-t-il ? » mais « ai-je déjà vu ce test
échouer ? ». C'est pour cette raison que les trois détecteurs de ce projet ont
été éprouvés en provoquant la panne qu'ils prétendent détecter.

\subsection{Perspectives}

Trois évolutions seraient utiles sans être nécessaires.

Le chiffrement du trafic suppose un nom de domaine : c'est la limite la plus
visible du déploiement actuel, et la plus simple à lever.

La haute disponibilité suppose de traiter l'écrivain unique autrement :
élection d'un responsable par document, ou externalisation de l'état. C'est un
chantier d'architecture à part entière, pas un réglage.

Le palier de rédacteurs, enfin, se mesure avant de s'optimiser : le résultat
réel doit guider la suite, plutôt qu'une optimisation supposée.

\subsection{Mot de la fin}

Le sujet prévenait : une application simple dont chaque choix est compris vaut
mieux qu'une application complexe qu'on ne sait pas expliquer. Ce dossier a
suivi cette consigne à la lettre, y compris dans ce qu'elle a de désagréable :
les limites sont énoncées, les hypothèses sont marquées comme telles, et les
écarts entre le plan et l'implémentation sont rapportés plutôt que lissés. C'est
la partie du travail qui a le plus servi, parce que c'est celle qui a révélé les
défauts que les tests ne voyaient pas.
